\section{代理制度}
\begin{enumerate}
  \item 基本原理。
        \begin{enumerate}
          \item 范围。\zd{仅限民事法律行为}（主要是合同）。\zd{事实行为、需本人亲自实施的不适用代理}。\al{离婚}。
          \item 类型。以\zd{谁的名义}实施的民事法律行为。\al{甲委托乙去云南购买普洱茶}。
                \begin{enumerate}
                  \item 直接代理。\al{乙以甲名义和丙签合同}。
                  \item 间接代理。\al{乙以自己名义}。\yi 显名（没啥好考的）。\al{丙知道甲}。\er 隐名。\al{丙不知道甲。乙没付钱}。
                        \begin{enumerate}
                          \item 代理人\zd{披露}义务。利益承受理论 + 合同相对性 \yjt\, \al{丙}择一选择（简单形成权）。
                          \item 被代理人\zd{介入权}。\al{甲给了钱，没收到茶叶}。
                        \end{enumerate}
                \end{enumerate}
          \item 代理关系的处理。
                \begin{enumerate}
                  \item 授权委托书。姓名、代理事项、权限、期限、\zd{被代理人签章}。
                  \item \zd{无因行为}。基础法律关系和\zd{代理权授予行为}不存在必然的因果关系 \yjt \zd{内外有别}。\al{票据、授权行为}。
                  \item 委托（2 方、合同）、代理（\zd{3 方。被代理人、代理人、第三人}）、委托代理不是一回事。
                        \begin{enumerate}
                          \item 有委托不一定有代理。\al{委托监护}。
                          \item 有代理不一定有委托。\al{法定监护}。
                        \end{enumerate}
                \end{enumerate}
          \item \zd{代理权的滥用}。\zd{是有权代理}！当代理由 3 方变 2 方的时候容易出现道德风险。
                \begin{enumerate}
                  \item 自己代理。一身二职 \yjt \zd{相对无效}。\al{甲委托乙买茶叶，乙将自家茶叶卖给甲}。
                  \item 双方代理。一仆二主 \yjt \zd{相对无效}。\al{甲委托乙买茶叶，丙委托乙卖茶叶，乙将丙茶叶卖给甲}。
                  \item \zd{恶意代理}。\zd{合同无效、共同侵权、连带责任、择一选择}。考的最多。\al{吃里扒外、恶意串通}。
                \end{enumerate}
        \end{enumerate}
  \item 分类。
        \begin{enumerate}
          \item 法定代理。\al{法定监护、甲律师接受法院指定担任被告人乙的辩护人}。
          \item 委托代理/意定代理。\zd{强烈的人身信赖关系}。
          \item 复代理/再代理/转委托。\al{甲特级茶 \yjt 乙忘告知 \yjt 丙胁迫 \yjt 丁茶农}。
                \begin{enumerate}
                  \item 复代理人。\zd{代理人自己名义选}。一般仅就\zd{第三人的选任、指示}承担过错责任。
                  \item 事前授权 \hz 同意、事后追认、代理人\zd{紧急情况为被代理人利益}。
                \end{enumerate}
          \item 职务代理/商事代理。
                \begin{enumerate}
                  \item 被代理人非自然人，是组织。
                  \item 本职工作内、\zd{内外有别}。
                \end{enumerate}
        \end{enumerate}
  \item 共同代理。多个代理人侵权问题。
  \item \zd{无权代理}。\zd{没有代理权、超越代理权、代理权过期}。
        \begin{enumerate}
          \item 效力待定。\zd{绝对不能认为是无效和可撤销}！\al{\xhx{伪造}公章、擅自、\xhx{冒用}}。
          \item 善意相对人 \yjt 催告权、\zd{选择权}（请求无权代理人履行 \hz 赔偿）。
          \item 被代理人 \yjt 否认权。核心考点。
          \item 无权代理人\zd{是否构成不当得利}，得先看被代理人态度，\zd{不能一概而论}。
          \item \bdy 无权代表！\al{公司董事长越权以公司名义为他人提供担保}。
        \end{enumerate}
  \item 表见代理/外观授权代理。\zd{相对人有理由相信无权代理人有代理权}（保护相对人信赖利益）。外观主义。
        \begin{enumerate}
          \item 主客观相结合。
                \begin{enumerate}
                  \item 主观。善意（\al{不知情、不知悉}） + 无过失（\zd{形式审查}权利外观。\al{核实 ×}）。
                  \item 客观。审查权利外观。\al{授权委托书、介绍信、交易习惯、\xhx{盖有公章 \hz 合同专用章}的空白合同书}。
                \end{enumerate}
          \item \zd{无权代理是常态，表见代理是例外}。十有八九构不成表见代理，慎选。
        \end{enumerate}
  \item vs 居间合同（仅仅是\zd{提供}媒介\zd{服务}）。
\end{enumerate}
